\scanpage{111}
\uline{Question 2} Soit $f:X\mathbin{-}S$ un morphisme \deleted{\unsure{affine}} de type fini de schémas affines noethériens, $f$ \deleted{\unsure{est plat local}} une intersection complète relative de dim relative $\geq3$ \deleted{p.ex.} p.ex. $f$ simple, et supposons $S$ normal \inserted{irréductible}. Soit $D$ sur $X$ satisfaisant les conditions a), b) du th. Alors est-il vrai encore que $D$ est un diviseur de Cartier ??

\figure{\begin{tikzpicture}\draw (0,0)--(.9,0);\draw (.13,-.12)--(.8,-.12);\end{tikzpicture}}

\uline{Corollaire du th.} Soit $R$ graphe d'équivalence propre dans un $S$-préschéma $X$, on suppose que $R\longrightarrow X$ est \uline{plat} \uline{propre} à fibres \inserted{géométriques} $(S_2)$ et \uline{intègres}, On suppose enfin $X$ \uline{normal} \inserted{intègre}, \deleted{\unsure{plus}} [à substituer à l'hyp. $X/R$ normal, $X/R$ \unsure{n'existant pas}\ldots]. Soit $D$ un \deleted{diviseur} \inserted{$1$-cycle} sur $X$,

\margin{et tel que $\underline{\operatorname{Pic}}_{R^{(1)}/X}$ existe et \unsure{soit} propre sur $\unsure{S|X}$}
\margin{\diagram{\begin{tikzcd}[row sep=large]
 R\arrow[r,shift left=.6ex]\arrow[r,shift right=.6ex] & X\\
 & S
\end{tikzcd}}}
\margin{\figure{\begin{tikzpicture}\draw (0,0)--(2.6,.08);\draw (2.7,1.8)--(2.72,-2.1);\end{tikzpicture}}}
\margin{\diagram{\begin{tikzcd}[row sep=small,column sep=small]
 D & D'=D_R^{(1)} & D'_\eta\\
 X\arrow[d,dash] & R\arrow[l,"p_1"']\arrow[d,"p_2"]\arrow[r] & R_\eta^{(2)}\arrow[d]\\
 {}\arrow[r,dash] & X & \eta
\end{tikzcd}}}
\margin{110}
